% geometry/curvature_space.tex

\chapter{Curvature Space}

\label{chap:curvature-space}

% ============================================================
\section*{Motivation}
% ============================================================

Classical railway alignment theory is usually formulated in position
space.

Within AIM, the primary geometric object is not the position curve
itself, but the curvature evolution along arc-length.

This shifts the interpretation of railway geometry from coordinate-based
description toward intrinsic curvature structure.

The concept of curvature space provides the mathematical setting for:
\begin{itemize}
\item transition curves,
\item sparse reconstruction,
\item realization analysis,
\item runtime evaluation,
\item and alignment optimization.
\end{itemize}

\begin{principle}[Curvature-space principle]
\label{princ:curvature-space}
Within AIM, railway alignment geometry is interpreted primarily as a
problem in curvature space.
\end{principle}

% ============================================================
\section{Curvature as Primary Geometry}
% ============================================================

\begin{definition}[Curvature law]
\label{def:curvature-law}
A curvature law is a function:
\[
\kappa : I \to \mathbb R,
\]
defined on an arc-length interval \(I\subset\mathbb R\).
\end{definition}

Within AIM, curvature is interpreted as the primary intrinsic geometric
quantity.

Geometry is reconstructed from curvature through integration.

\begin{principle}[Curvature-first principle]
\label{princ:curvature-first}
Within AIM, geometry is generated from curvature rather than curvature
being derived from geometry.
\end{principle}

% ============================================================
\section{Curvature Space}
% ============================================================

\begin{definition}[Curvature space]
\label{def:curvature-space}
Curvature space denotes the space of admissible curvature laws:
\[
\mathcal K
=
\{
\kappa : I \to \mathbb R
\}.
\]
\end{definition}

The symbol \(\mathcal K\) does not denote one fixed function class.
It denotes the conceptual space in which different admissibility,
smoothness, realization, and optimization requirements may be imposed.

Different subclasses of curvature space correspond to different
geometric and operational properties.

Examples include:
\begin{itemize}
\item constant curvature,
\item piecewise constant curvature,
\item continuous curvature,
\item differentiable curvature,
\item bounded curvature,
\item transition curvature,
\item and hybrid curvature laws.
\end{itemize}

\OPEN{
Later extensions of curvature space include admissible curvature laws,
curvature spectra, realization-stable curvature, curvature energy,
curvature bandwidth, transition-family subspaces, Vienna-type
reinterpretations, and operator deformations.
}

% ============================================================
\section{Intrinsic Reconstruction}
% ============================================================

\begin{definition}[Intrinsic reconstruction]
\label{def:intrinsic-curvature-reconstruction}
Given a curvature law:
\[
\kappa(s),
\]
the tangent direction satisfies:
\[
\theta'(s)=\kappa(s),
\]
and geometry is reconstructed through:
\[
\gamma'(s)=\mathbf t(s).
\]
\end{definition}

Curvature therefore acts as the differential generator of geometry.

The position curve \(\gamma\) is not primary. It is the result of
integrating the intrinsic curvature state together with initial pose
data.

% ============================================================
\section{Constant Curvature}
% ============================================================

\begin{definition}[Constant curvature]
\label{def:constant-curvature}
A constant curvature law satisfies:
\[
\kappa(s)=\kappa_0.
\]
\end{definition}

Special cases are:
\begin{itemize}
\item \(\kappa_0=0\): straight line,
\item \(\kappa_0\neq 0\): circular arc.
\end{itemize}

Constant-curvature elements form the fixed parts of the sparse alignment
model.

% ============================================================
\section{Transition Curvature}
% ============================================================

\begin{definition}[Transition curvature]
\label{def:transition-curvature}
A transition curvature law connects different curvature states
continuously or smoothly.
\end{definition}

Transition curves are therefore interpreted as curvature-transition
processes rather than merely geometric interpolation curves.

Typical transition quantities include:
\[
\kappa'(s),
\qquad
\kappa''(s),
\qquad
\kappa'''(s).
\]

In AIM, transition families are interpreted as structured subspaces or
parameterized paths within curvature space.

% ============================================================
\section{Curvature Continuity}
% ============================================================

\begin{definition}[Curvature continuity]
\label{def:curvature-continuity}
A curvature law is:
\begin{itemize}
\item \(C^0\)-continuous if \(\kappa(s)\) is continuous,
\item \(C^1\)-continuous if \(\kappa'(s)\) is continuous,
\item \(C^2\)-continuous if \(\kappa''(s)\) is continuous.
\end{itemize}
\end{definition}

Higher continuity improves:
\begin{itemize}
\item dynamic smoothness,
\item realization stability,
\item numerical robustness,
\item and passenger comfort.
\end{itemize}

Continuity of position is not sufficient for high-quality railway
geometry. The decisive structure lies in the continuity and behaviour of
curvature and its derivatives.

% ============================================================
\section{Sparse Curvature Representation}
% ============================================================

\begin{definition}[Sparse curvature representation]
\label{def:sparse-curvature-representation}
A sparse representation describes curvature using:
\begin{itemize}
\item fixed curvature elements,
\item transition elements,
\item and compact parameter sets.
\end{itemize}
\end{definition}

Within AIM, sparse alignments are interpreted as compact curvature
programs.

The sparse model stores:
\begin{itemize}
\item curvature structure,
\item transition families,
\item and reconstruction rules,
\end{itemize}
rather than dense geometric point clouds.

This makes sparse alignment suitable for:
\begin{itemize}
\item reconstruction,
\item import normalization,
\item runtime evaluation,
\item realization comparison,
\item and optimization.
\end{itemize}

% ============================================================
\section{Hybrid Curvature Representation}
% ============================================================

\begin{definition}[Hybrid curvature model]
\label{def:hybrid-curvature-model}
A hybrid curvature model combines:
\begin{itemize}
\item parametric curvature structure,
\item sampled correction data,
\item and local refinement.
\end{itemize}
\end{definition}

Hybrid representations allow:
\begin{itemize}
\item compact storage,
\item efficient reconstruction,
\item realization-aware correction,
\item measured-state integration,
\item and scalable runtime evaluation.
\end{itemize}

Hybrid curvature models are especially relevant when measured or
realized geometry cannot be represented adequately by a purely analytical
target model.

% ============================================================
\section{Normalized Curvature Space}
% ============================================================

\begin{definition}[Normalized curvature law]
\label{def:normalized-curvature-law}
A normalized curvature law is defined on:
\[
w\in[0,1].
\]
\end{definition}

Normalized representations separate:
\begin{itemize}
\item geometric shape,
\item from physical scaling.
\end{itemize}

This enables:
\begin{itemize}
\item reusable transition families,
\item lookup-based transitions,
\item family interpolation,
\item and normalized optimization.
\end{itemize}

Physical curvature laws are obtained from normalized laws by scaling in
length and curvature amplitude.

% ============================================================
\section{Curvature Derivatives}
% ============================================================

\begin{definition}[Curvature derivatives]
\label{def:curvature-derivatives}
Important higher-order quantities include:
\[
\kappa'(s),
\qquad
\kappa''(s),
\qquad
\kappa'''(s).
\]
\end{definition}

These quantities influence:
\begin{itemize}
\item jerk,
\item realization behaviour,
\item dynamic smoothness,
\item vehicle response,
\item and transition quality.
\end{itemize}

Curvature derivatives provide the bridge from intrinsic geometry to
runtime dynamics.

% ============================================================
\section{Frequency Interpretation}
% ============================================================

Curvature laws may also be interpreted spectrally.

Low-frequency curvature components determine global geometric structure,
whereas high-frequency components influence:
\begin{itemize}
\item realization sensitivity,
\item numerical instability,
\item measurement sensitivity,
\item and dynamic roughness.
\end{itemize}

\begin{proposition}[Spectral realization principle]
\label{prop:spectral-realization-principle}
Realization acts approximately as a low-pass filter on curvature space.
\end{proposition}

This spectral view connects curvature space directly to construction,
maintenance, and physical realization.

% ============================================================
\section{Curvature Space and Runtime}
% ============================================================

Runtime evaluation operates primarily on curvature structures rather
than directly on dense geometry.

Projection, reconstruction, frame evaluation, pose3 evaluation, and
dynamics are derived from runtime evaluation of curvature-based states.

Thus, curvature space is not only a design space. It is also the
computational state space underlying runtime reconstruction.

% ============================================================
\section{Curvature Space and Realization}
% ============================================================

Realization transforms target curvature into physically realized
curvature.

This may be written abstractly as:
\[
\kappa_{\mathrm{real}}
=
\mathcal R_{\kappa}
(
\kappa_{\mathrm{target}}
).
\]

The realization operator may attenuate, deform, or locally modify
curvature evolution.

Consequently, the physically relevant design object is not only the
target curvature law, but also its expected realized curvature state.

% ============================================================
\section{Curvature Space and Optimization}
% ============================================================

Optimization within AIM primarily acts in curvature space.

Typical optimization variables include:
\begin{itemize}
\item transition parameters,
\item curvature amplitudes,
\item transition lengths,
\item curvature anchors,
\item and sparse curvature corrections.
\end{itemize}

AIM optimization may therefore be interpreted as the search for
admissible, realizable, and operationally suitable paths in curvature
space.

% ============================================================
\section{Canonical Interpretation}
% ============================================================

\begin{proposition}
Within AIM, railway alignment geometry is fundamentally a curvature-space
problem.
\end{proposition}

\begin{proposition}
Position-space geometry is interpreted as reconstructed geometry derived
from curvature evolution.
\end{proposition}

\begin{proposition}
Transition curves are interpreted as controlled transitions within
curvature space.
\end{proposition}

\begin{proposition}
Realization, runtime evaluation, and optimization act on curvature-space
structures before they act on visualized position geometry.
\end{proposition}

% ============================================================
\section{Connection to AIM}
% ============================================================

\begin{summary}
Curvature space forms the intrinsic geometric foundation of AIM.

Within this framework:
\begin{itemize}
\item curvature acts as primary geometry,
\item sparse models become curvature programs,
\item transition curves become curvature transitions,
\item runtime systems operate on intrinsic curvature structure,
\item realization transforms curvature states,
\item and optimization searches within curvature-space structures.
\end{itemize}

This interpretation unifies geometry, realization, runtime evaluation,
measurement interpretation, and optimization within one intrinsic
mathematical framework.
\end{summary}
